%=====================================
% Khai báo nhóm Tex (cơ bản) - MINIMALIST VERSION + WATERMARK + FOOTER
%=====================================
%===========Bảng
\usepackage{longtable,multirow,makecell,tabularx}
\usepackage{diagbox}
\usepackage[shortlabels]{enumitem}
% xcolor đã được load trong mausac-minimal.tex
%---
\let\Bbbk\relax
\usepackage{extarrows,amsmath,amssymb,mathrsfs,maybemath,xlop,polynom,slashbox}
%\usepackage{yhmath} % Bỏ để tránh xung đột \widering với fouriernc
%\usepackage{enumerate}
\usepackage{tikz} 
\usepackage{tkz-euclide}
\usepackage{tikz-3dplot}
\usepackage{tkz-tab}
\usepackage{pifont}
\usepackage{bbding}
\usepackage{stackengine}
\usepackage{scalerel}
\usepackage{array}
\usepackage{xcolor,colortbl,tasks}
%==========
\usetikzlibrary{math,through,calc,intersections,angles,quotes,shapes,shapes.geometric,arrows,patterns,snakes,matrix,chains,arrows.meta,decorations.shapes,decorations.fractals,decorations.markings,shadows}
\usetikzlibrary{positioning,decorations.text,decorations.pathmorphing}
\usetikzlibrary{shadings,fadings}
\usepackage{pgfplots}
\usepackage{pgfornament}
\usepgfplotslibrary{fillbetween}
\pgfplotsset{compat=1.9}
\usepackage[hidelinks,unicode]{hyperref}
\usepackage{pdfpages} % Cho phép chèn file PDF với \includepdf
\usepackage{currfile}
\usepackage[outline]{contour}
\usepackage{fontawesome}
\usepackage{lipsum}
\usepackage{everypage,eso-pic}
\usepackage{microtype} % Cải thiện khoảng cách giữa các từ/ký tự
\setlength{\emergencystretch}{3em} % Giúp giãn dòng đều hơn, tránh lỗi underfull

\arrayrulecolor{brandNavy}
\newcolumntype{C}[1]{>{\centering\arraybackslash}p{#1}}
\newcolumntype{L}[1]{>{\raggedright\arraybackslash}p{#1}}

\usepackage[top=\tren cm, bottom=\duoi cm, left=\trai cm, right=\phai cm, headheight=24pt] {geometry}
\usepackage{varwidth}

% Thiết lập cho layout 2 cột
\setlength{\columnsep}{8mm} % Khoảng cách giữa 2 cột
\setlength{\columnseprule}{0.4pt} % Đường kẻ mảnh giữa 2 cột
\def\columnseprulecolor{\color{warmgray}} % Màu đường kẻ
%--------------Gói trắc nghiệm EX-TEST
\usepackage[loigiai]{ex_test} 
\ifdefined\BTRL\else\def\BTRL{}\fi 

\font\damEX=ugqb8v at 13pt
\def\loigiaiEX{\makebox[\linewidth][c]{\rule[-0.2cm]{0pt}{0.6cm}\color{brandNavy}\damEX\faPencilSquareO\ Lời giải}}
\renewcommand{\nameex}{\damEX\color{\mauEX} \iconVD Câu}
\def\qedEX{\color{\mauDC}\ensuremath{\square}}
\newcommand{\hoac}[1]{\left[\begin{aligned}#1\end{aligned}\right.}
\newcommand{\heva}[1]{\left\{\begin{aligned}#1\end{aligned}\right.}

%----------
\usepackage{esvect}
\def\vec{\vv}
\def\overrightarrow{\vv}
\DeclareSymbolFont{symbolsC}{U}{txsyc}{m}{n}
\DeclareMathSymbol{\varparallel}{\mathrel}{symbolsC}{9}
\DeclareMathSymbol{\parallel}{\mathrel}{symbolsC}{9}

%=====================================
% TÊN DỌC BÊN LỀ PHẢI - VINHMATH WATERMARK VỚI RUY BĂNG
%=====================================
\AddToShipoutPictureBG{%
\ifnum\value{page}>0
\begin{tikzpicture}[remember picture,overlay]
% Faded logo watermark in the center of the page
\node[opacity=0.08] at (current page.center) {
    \includegraphics[width=10.5cm]{logo/CLB-MAP-logo-circle-4K.png}
};
\coordinate (margin) at ([xshift=-5mm]current page.east);
% Ruy băng phía trên
\draw[brandBlue!30,line width=0.8pt] ([yshift=3.5cm]margin) -- ([yshift=7cm]margin);
\fill[brandGold] ([yshift=7.2cm]margin) circle (1.5pt);
% Tên ở giữa
\node[rotate=90,anchor=center,font=\fontfamily{qag}\fontseries{b}\fontsize{10}{12}\selectfont,text=brandGold!45] at (margin) {\watermarktext};
% Ruy băng phía dưới
\draw[brandBlue!30,line width=0.8pt] ([yshift=-3.5cm]margin) -- ([yshift=-7cm]margin);
\fill[brandGold] ([yshift=-7.2cm]margin) circle (1.5pt);
\end{tikzpicture}%
\fi
}

%--------------------------
% HEADER AND FOOTER STYLING - MINIMALIST TWOSIDE
%--------------------------
\usepackage{fancyhdr,lastpage}
\pagestyle{fancy}
\fancyhf{}
\renewcommand{\headrulewidth}{0.5pt}
\renewcommand{\footrulewidth}{0.4pt}
\renewcommand{\headrule}{\hbox to\headwidth{\color{\maufoot}\leaders\hrule height \headrulewidth\hfill}}
\renewcommand{\footrule}{\hbox to\headwidth{\color{\maufoot!30}\leaders\hrule height \footrulewidth\hfill}}

% Header thiết kế đối xứng (Twoside)
\fancyhead[LE]{\color{\maufoot}\sffamily\bfseries\leftmark}
\fancyhead[RO]{\color{\maufoot}\sffamily\bfseries\rightmark}
\fancyhead[RE,LO]{\color{\maufoot}\sffamily\fontsize{9pt}{11pt}\selectfont\itshape M.A.P - Math And Passion}

% Footer thiết kế đối xứng (Twoside)
\fancyfoot[LE,RO]{\color{\maufoot}\fontsize{9pt}{11pt}\selectfont\sffamily Trang \thepage/\pageref{LastPage}}
\fancyfoot[RE,LO]{\color{warmgray}\fontsize{9pt}{11pt}\selectfont\faPhone\ \SDT}

% Định nghĩa lại style plain cho trang đầu chương (đối xứng)
\fancypagestyle{plain}{%
    \fancyhf{}%
    \renewcommand{\headrulewidth}{0pt}%
    \renewcommand{\footrulewidth}{0.4pt}%
    \renewcommand{\footrule}{\hbox to\headwidth{\color{\maufoot!30}\leaders\hrule height \footrulewidth\hfill}}%
    \fancyfoot[LE,RO]{\color{\maufoot}\fontsize{9pt}{11pt}\selectfont\sffamily Trang \thepage/\pageref{LastPage}}%
    \fancyfoot[RE,LO]{\color{\maufoot}\fontsize{8.5pt}{11pt}\selectfont\sffamily\bfseries\itshape M.A.P - Math And Passion}%
}


% Logo commands (kept for compatibility but simplified output if needed)
\def\hienLOGO{} 
\def\anLOGO{}

%--------------------------
% SECTION STYLING - MINIMALIST
%--------------------------
\usepackage[explicit]{titlesec}
\usepackage{titledot}

\setcounter{secnumdepth}{4}
\renewcommand{\thechapter}{\arabic{chapter}}
\renewcommand{\thesection}{\arabic{section}}
\renewcommand{\thesubsection}{\Alph{subsection}}
\renewcommand{\thesubsubsection}{\arabic{subsubsection}}

% Chapter
\titleformat{\chapter}[display]
  {\color{\mauchuong}\fontfamily{qag}\bfseries\Huge}
  {\flushright\fontsize{50}{60}\selectfont\thechapter}
  {10pt}
  {\titlerule[2pt]\vspace{1ex}\flushright\Huge\MakeUppercase{#1}}
\titlespacing*{\chapter}{0pt}{-40pt}{20pt}

% Section
\titleformat{\section}
  {\color{\mauSEC}\fontfamily{qag}\bfseries\Large}
  {\thesection.}
  {0.5em}
  {#1}
  [\titleline{\color{\mauSEC}\titlerule[1pt]}]

% Subsection
\titleformat{\subsection}
  {\color{\mauSUBSEC}\fontfamily{qag}\bfseries\large}
  {\thesubsection.}
  {0.5em}
  {#1}

% Subsubsection
\titleformat{\subsubsection}
  {\color{\mauSUBSUBSEC}\fontfamily{qag}\bfseries\normalsize}
  {\thesubsubsection.}
  {0.5em}
  {#1}

% Boxes and Environments - MINIMALIST
\usepackage[most]{tcolorbox}

% Define simple clean boxes
\def\beginboxVD{
\begin{tcolorbox}[
    enhanced,
    breakable,
    colback=\mauVD!5,
    colframe=\mauVD,
    coltitle=\mauVD,
    boxrule=0pt,
    leftrule=3pt,
    sharp corners,
    left=2mm,right=2mm,top=2mm,bottom=2mm,
    before skip=2mm, after skip=2mm,
    parbox=false
]
}

\def\beginboxEX{
\begin{tcolorbox}[
    enhanced,
    breakable,
    colback=white,
    colframe=\mauEX,
    boxrule=0pt,
    leftrule=3pt,
    sharp corners,
    left=2mm,right=2mm,top=2mm,bottom=2mm,
    before skip=2mm, after skip=2mm,
    parbox=false
]
}

\def\beginboxBT{
\begin{tcolorbox}[
    enhanced,
    breakable,
    colback=white,
    colframe=\mauBT,
    boxrule=0pt,
    leftrule=3pt,
    sharp corners,
    left=2mm,right=2mm,top=2mm,bottom=2mm,
    before skip=2mm, after skip=2mm,
    parbox=false
]
}

\def\endbox{\end{tcolorbox}}

%=====================================
% Counter cho Dạng toán
%=====================================
\newcounter{dang}\setcounter{dang}{0}
\renewcommand{\thedang}{\arabic{dang}}
\makeatletter
\@addtoreset{dang}{subsection}
\@addtoreset{bt}{subsection}
\@addtoreset{ex}{subsection}
\@addtoreset{vd}{subsection}
\makeatother

% DANG (Type of problem) env
\newtcolorbox{dangg}[1]{
    enhanced,
    breakable,
    colback=white,
    colframe=\maudang,
    coltitle=white,
    colbacktitle=\maudang,
    boxrule=0pt,
    toprule=2pt,
    sharp corners,
    fonttitle=\bfseries\sffamily,
    title={Dạng \stepcounter{dang}\thedang. #1},
    attach boxed title to top left={yshift=-2mm, xshift=0mm},
    boxed title style={sharp corners, boxrule=0pt},
    left=2mm,right=2mm,top=4mm,bottom=2mm,
    before upper={\setlength{\parskip}{0.3em}}
}
\newenvironment{dang}[1]{\begin{dangg}{#1}}{\end{dangg}}

%=====================================
% Môi trường Định nghĩa, Định lí
%=====================================
\newenvironment{boxdl}
{\begin{tcolorbox}
		[enhanced jigsaw,breakable,pad at break*=1mm,
		colback=\nendl,boxrule=0pt,frame hidden, 
		opacityback=0.5,
		left=2mm, right=0pt, bottom=1.5pt, top=1.5pt,
		before skip=2mm,
		after skip=2mm,
		overlay={\draw[double,line width=1.5pt,\maudl] ([xshift=3pt]interior.north west)--([xshift=3pt]interior.south west);}
		]}
	{\end{tcolorbox}}

\newenvironment{boxdn}
{\begin{tcolorbox}
		[enhanced jigsaw,breakable,pad at break*=1mm,
		colback=\maudn!3,boxrule=0pt,frame hidden, 
		opacityback=0.5,
		left=2mm, right=0pt, bottom=1.5pt, top=1.5pt,
		before skip=2mm,
		after skip=2mm,
		overlay={\draw[double,line width=1.5pt,\maudn] ([xshift=3pt]interior.north west)--([xshift=3pt]interior.south west);}
		]}
	{\end{tcolorbox}}

\newenvironment{dn}{\begin{boxdn}\begin{mydn}}{\end{mydn}\end{boxdn}}
\newenvironment{dl}{\begin{boxdl}\begin{mydl}}{\end{mydl}\end{boxdl}}

%=====================================
% Môi trường Tóm tắt
%=====================================
\newenvironment{tomtat}{}{}

%=====================================
% Môi trường Chú ý và Lưu ý
%=====================================
\newenvironment{note}
{\begin{tcolorbox}
		[enhanced jigsaw,breakable,pad at break*=1mm,
		opacityback=0,boxrule=0pt,frame hidden,
		left=8mm, right=0pt, bottom=0pt, top=0pt,
		before skip=1mm,
		after skip=1mm,
		underlay unbroken and first={
			\node[\maucy,font=\large\bfseries] at ([xshift=0.3cm,yshift=-0.32cm]interior.north west) {\faExclamationTriangle};
		},
		fontupper=\it,
		]}
	{\end{tcolorbox}}

\newenvironment{luuy}
{\begin{tcolorbox}
		[enhanced jigsaw,breakable,pad at break*=1mm,
		opacityback=0,boxrule=0pt,frame hidden,
		left=8mm, right=0pt, bottom=0pt, top=0pt,
		before skip=1mm,
		after skip=1mm,
		underlay unbroken and first={
			\node[\maucy,font=\large\bfseries] at ([xshift=0.3cm,yshift=-0.32cm]interior.north west) {\faExclamationTriangle};
		},
		fontupper=\it,
		]}
	{\end{tcolorbox}}

% Other environments
\newtheorem{mydn}{\color{\maudn}\damEX Định nghĩa}[section]
\newtheorem{vd}{\color{\mauVD}\damEX\iconVD Ví dụ}
\newtheorem{bt}{\color{\mauBT}\damEX\iconVD Bài}
\newtheorem{ap}{\color{\mauBT}\damEX\iconVD Áp dụng}
\newtheorem{myvidu}{\color{\mauVD}\damEX\iconVD Ví dụ}
\newtheorem{mydl}{\color{\maudl}\damEX Định lí}
\newtheorem{nx}{\color{\maucy}\damEX Nhận xét}
\newtheorem{phantich}{\color{\maupt}\damEX Phân tích}
\newtheorem{tc}{\color{\maudl}\damEX Tính chất}

% Helper macros
\def\khoanhtrondapan{}
\def\khongkhoanhtrondapan{}
\def\FULLWIDTH{\newgeometry{top=\tren cm, bottom=\duoi cm, left=\trai cm, right=\phai cm, headheight=24pt}}
\def\NOTE{\newgeometry{top=\tren cm, bottom=\duoi cm, left=\trai cm, right=\phai cm, headheight=24pt}}
\def\indapan#1#2{}
% Các biến này được định nghĩa trong file main, không định nghĩa lại ở đây
% \def\chamngon{}
% \def\trenphai{}
% \def\duoiphai{}
% \def\hoten{}
% \def\SDT{}
% \def\DC{}

% Canh chỉnh mục lục chính
\setcounter{secnumdepth}{4}
\setcounter{tocdepth}{3}
\contentsmargin{1.5cm}

% Vẽ số tròn cho Subsection trong mục lục
\newcommand{\tron}[1]{%
	\begin{tikzpicture}[baseline=(A.base)]%
		\node[circle,draw=brandGold,line width=0.5pt,fill=white,inner sep=2pt,outer sep=1.5pt] (A) {\color{brandGold} #1};
		\node[circle,draw=white,fill=brandGold,inner sep=1pt,outer sep=1pt] (A) {\bfseries\color{white} #1};
	\end{tikzpicture}%
}

\titlecontents{part}[0pc]
{\addvspace{12pt}\color{brandNavy}\fontsize{16pt}{18pt}\selectfont\bfseries}
{Phần \thecontentslabel\quad}
{}
{\hfill\contentspage}
[\addvspace{4pt}{\color{brandGold}\titlerule[1.5pt]}\addvspace{8pt}]

\titlecontents{chapter}[5.2pc]
{\addvspace{8pt}\color{brandNavy}\fontsize{11pt}{13pt}\selectfont\bfseries\sffamily}
{\contentslabel[\chaptername\ \thecontentslabel.]{5.2pc}}
{}
{\dotfill\contentspage}

\titlecontents{section}[6.5pc]
{\addvspace{4pt}\color{brandBlue}\fontsize{10.5pt}{12pt}\selectfont\bfseries\sffamily}
{\contentslabel[\thecontentslabel.]{1.3pc}}
{}
{{\scriptsize\dotfill}\contentspage}

\titlecontents{subsection}[8.0pc]
{\addvspace{2pt}\color{brandGold}\fontsize{10pt}{11.5pt}\selectfont\sffamily\bfseries}
{\contentslabel[\tron{\thecontentslabel}]{1.5pc}}
{}
{{\tiny\dotfill}\contentspage}

\titlecontents{subsubsection}[9.5pc]
{\addvspace{1pt}\color{brandDark}\fontsize{9.5pt}{11pt}\selectfont\itshape}
{\contentslabel[\thecontentslabel]{1.5pc}}
{}
{{\tiny\dotfill}\contentspage}

% TOC
\def\muclucchinh{
    \tableofcontents
    \pagenumbering{roman}
}

% Fix for ex_test compatibility
\def\dotlineEX#1{\par\noindent\dotfill\par}
\def\mauDC{\mauLG!50}
\setlength{\parindent}{0pt}
